% Student assessment with shared general tasks and criterion-specific Available points tables; no solutions.
% Scoring ledger: Q1 and Q2 C1=2+2=4, C2=1+2+1=4 (8 each).
% Q3 C1=3+3=6, C2=2+2+2=6 (12); totals C1=14, C2=14, overall=28.
% Scenario validity is assessed within Q3 task marks; no construction bonus.
% Build twice with ../../../preamble/ in TEXINPUTS.
\makeatletter
\def\input@path{{../../../preamble/}}
\makeatother
\input{preamble_exams}
\graphicspath{{../../../preamble/}}
\usepackage{needspace}
\hypersetup{pdftitle={Grade 11 Term 2 Learning Evidence 1: C1 and C2},pdfauthor={Global Economics}}
\setYColumnWidth{0.36}
\linespread{1.04}
\raggedbottom
\newcolumntype{K}{>{\raggedright\arraybackslash}X}
\newcommand{\examtable}[2]{\par\noindent\begin{minipage}{\linewidth}\small
\renewcommand{\arraystretch}{1.08}\setlength{\tabcolsep}{3pt}
\begin{tabularx}{\linewidth}{#1}\toprule #2\bottomrule\end{tabularx}
\end{minipage}\par}
\let\ExamSubsectionBox\SubsectionBox
\RenewDocumentCommand{\SubsectionBox}{s m}{%
  \par\needspace{4\baselineskip}% keep heading with opening prompt in multicols
  \begingroup
  \setlength{\linewidth}{\dimexpr\linewidth-2\fboxsep\relax}%
  \IfBooleanTF{#1}{\ExamSubsectionBox*{#2}}{\ExamSubsectionBox{#2}}%
  \endgroup\nobreak
}
% Shared headings use stepcounter, not refstepcounter: capture the automatic
% number explicitly for score-table references, without manual question IDs.
\makeatletter
\newcommand{\QuestionLabel}[1]{%
  \begingroup\edef\@currentlabel{\arabic{subsectionboxproblem}}%
  \label{#1}\endgroup
}
\makeatother
\newcommand{\cOneA}{Identify the population and both samples: place, period and sizes. Pair and interpret population and sample means, standard deviations and proportions.}
\newcommand{\cOneB}{Calculate the changes in sample mean, standard deviation, and proportion from the initial to the second independent study. Interpret these differences, identifying which measures changed and which remained unchanged.}
\newcommand{\cTwoA}{Complete the C2 information table for the initial study: initially unknown and later established mean, known spread, sample statistics and procedure settings.}
\newcommand{\cTwoB}{For each requested confidence level, tabulate the initial point estimate, margin of error, separate endpoints, confidence interval and full width using the known population standard deviation.}
\newcommand{\cTwoC}{For subsequent studies only, calculate covering counts, observed coverage and absolute differences from nominal confidence (percentage points).}

\begin{document}
\noindent
\begin{tabularx}{\textwidth}{|C{2.8cm}|C{\dimexpr\textwidth-5.6cm-6\tabcolsep-4\arrayrulewidth\relax}|C{2.8cm}|}
\hline
\centering\vspace{3mm}\includegraphics[width=2.5cm]{logo.png}
& \CellCenter{\textbf{GLOBAL ECONOMICS}\par
\textbf{Grade 11 -- Learning Evidence 1}\par
\textbf{C1 and C2 Integrated Exam}\par
\textbf{Teacher: Nicolás López Cuéllar}}
& \CellCenter{\textbf{TERM 2}\par\textbf{2026--2027}}\\\hline
\end{tabularx}\par\smallskip
\LearningObjective{Distinguish fixed population parameters from variable sample statistics in economic and financial contexts; calculate known-population-standard-deviation confidence intervals and interpret their long-run coverage.}{%
\textbf{C1}: Distinguishes population parameters from sample statistics by interpreting and comparing population and sample means, standard deviations and proportions in economic and financial contexts.\par
\textbf{C2}: Compares point and interval estimation and interprets confidence through long-run coverage, distinguishing this from probability assigned to a particular calculated interval containing the parameter.}
\par\smallskip
\StudentInfo\par

\begin{multicols}{2}
	
\setlength{\parskip}{3pt}
\SubsectionBox*{Criteria assessment}
To pass, students must obtain a \textbf{minimum of 4 points in each criterion}. Justify all answers when possible by showing relevant formulas, substitutions, calculations, and units. Partial credit may be awarded for correct procedures despite arithmetic errors; answers without supporting procedures receive no credit. A 1-point penalty per criterion applies if the student's name is missing or the assessment is submitted late.



\noindent\begin{minipage}{\linewidth}
\centering
\textbf{Available points}\par\smallskip

\centering
\begin{tabular}{|c||c|c||c|c|c|}
\hline
\multirow[c]{2}{*}{\textbf{\textnumero}} & \multicolumn{2}{c||}{\textbf{C1}}
& \multicolumn{3}{c|}{\textbf{C2}}\\
\cline{2-6}
& \textbf{T1} & \textbf{T2}
& \textbf{T1} & \textbf{T2} & \textbf{T3}\\
\hline
\ref*{q:1} & 2 & 2 & 1 & 2 & 1\\
\hline
\ref*{q:2} & 2 & 2 & 1 & 2 & 1\\
\hline
\ref*{q:3} & 3 & 3 & 2 & 2 & 2\\
\hline
\textbf{Total}
& \multicolumn{2}{c||}{14}
& \multicolumn{3}{c|}{14}\\
\hline
\end{tabular}

\end{minipage}\par


\raggedcolumns
\SubsectionBox*{General tasks for every problem}
All scenario data are hypothetical.

\begingroup\small
\textbf{C1 Task 1}. \cOneA\ Use two tables with these headings. For the population, include $\mathcal P,N,\mu,\sigma,p$; supply values known/unknown status, units and contextual meanings.

\vspace{-3mm}
\begin{center}
\begin{tabularx}{\linewidth}{@{}l | K@{}}
\toprule
\textbf{Symbol} & \textbf{Value and Meaning in context}\\
\bottomrule
\end{tabularx}
\end{center}

For the samples, include $X_A,X_B$, $n_A,n_B$, $\bar x_A,\bar x_B$, $s_A,s_B$, $k_A,k_B$ and $\hat p_A,\hat p_B$. Supply each values units and show both proportion calculations.
\begin{center}
\begin{tabularx}{\linewidth}{@{}l | K | K@{}}
\toprule
\textbf{Symbol} & \textbf{Initial sample A} & \textbf{Second sample B}\\
\bottomrule
\end{tabularx}
\end{center}

\columnbreak 

\textbf{C1 Task 2}. \cOneB\ Use rows for $\Delta\bar x$, $\Delta s$ and $\Delta\hat p$; show substitutions, signed results and units in the calculation column, and contextual comparisons in the interpretation column.
\begin{center}
\begin{tabularx}{\linewidth}{@{}l | K@{}}
\toprule
\textbf{Symbol} &  \textbf{Calculation, Value and Meaning in context} \\
\bottomrule
\end{tabularx}
\end{center}

\textbf{C2 Task 1}. \cTwoA\ Use contiguous groups in this order: Population parameters ($\mu,\mu_{\mathrm{known}},\sigma$); Sample statistics ($\bar x,n,\mathrm{SE}$); Procedure settings ($C=1-\alpha,\alpha,\alpha/2,z_{1-\alpha/2}$). Supply values or status, units and contextual meanings; show the SE calculation and procedure-setting calculations for every requested confidence level.
\begin{center}
\begin{tabularx}{\linewidth}{@{}l | K@{}}
\toprule
\textbf{Symbol} & \textbf{Value and Meaning in context}\\
\bottomrule
\end{tabularx}
\end{center}

\textbf{C2 Task 2}. \cTwoB\ Use a separate table for each requested confidence level, with rows $\bar x,\mathrm{MoE},L,U,\mathrm{CI},W$. Show calculations and then results with units and contextual meanings.
\begin{center}
\begin{tabularx}{\linewidth}{@{}l | K@{}}
\toprule
\textbf{Symbol} & \textbf{Calculation, Value and meaning in context}\\
\bottomrule
\end{tabularx}
\end{center}

\textbf{C2 Task 3}. \cTwoC\ Show calculations for $\mathrm{Cov}(C)$ and $E_{\mathrm{abs}}(C)$, with coverage in percent and absolute difference in percentage points. 

\endgroup

\columnbreak

\SubsectionBox*{Task reference material}
\begingroup\small
\textbf{C1 pairings.} Population $\mathcal P$, size $N$; sample $X$, size $n$; population/sample mean $(\mu,\bar x)$; standard deviation $(\sigma,s)$; proportion $(p,\hat p)$. A condition count is $k$, so $\hat p=k/n$. A sample-proportion change is $100(\hat p_B-\hat p_A)$ percentage points.

The derived SE is placed in the sample-statistics group; here it uses known $\sigma$ and fixed $n$, so it is not an estimate using $s$. The later-known notation refers to the same fixed mean, not a second parameter.

\textbf{Two-sided standard normal reference.} Use the following numerical critical values for all calculations.
\par\noindent\begin{minipage}{\linewidth}
\begin{minipage}[t]{.49\linewidth}\vspace{0pt}
\examtable{@{}KK@{}}{
\textbf{Confidence level}\newline $1-\alpha$&\textbf{Positive critical value}\newline $z_{1-\alpha/2}$\\\midrule
80\%&1.282\\
90\%&1.645\\
95\%&1.960\\
}
\end{minipage}\hfill
\begin{minipage}[t]{.49\linewidth}\vspace{0pt}
\examtable{@{}KK@{}}{
\textbf{Confidence level}\newline $1-\alpha$&\textbf{Positive critical value}\newline $z_{1-\alpha/2}$\\\midrule
98\%&2.326\\
99\%&2.576\\
99.9\%&3.291\\
}
\end{minipage}
\end{minipage}\par
% Retain the column break between the reference tables and formulas.

\[
\begin{aligned}
\alpha&=1-C, &\text{each tail}&=\alpha/2,\\
\mathrm{SE}&=\frac{\sigma}{\sqrt n}, &\mathrm{MoE}&=z_{1-\alpha/2}\mathrm{SE},\\
L&=\bar x-\mathrm{MoE}, &U&=\bar x+\mathrm{MoE},\\
\mathrm{CI}&=[L,U], &W&=U-L=2\mathrm{MoE}.
\end{aligned}
\]

For each confidence level separately, a subsequent interval covers when
$L_i\leq\mu_{\mathrm{known}}\leq U_i$. Equality at an endpoint counts as coverage.
\[
\mathrm{Cov}(C)=\frac{N_{\mathrm{cover}}(C)}{N(C)}\,100\%.
\]
\[
E_{\mathrm{abs}}(C)=|\text{observed percentage}-100C|
\]
measured in percentage points.
Proportions and confidence settings are unitless; a mean, standard deviation, standard error, margin, endpoint and width use the measurement's units.

\endgroup

\columnbreak
\begingroup\small
\SubsectionBox{Young adults' grocery budgets \score{C1-\sml{4}\quad C2-\sml{4}}}
\QuestionLabel{q:1}

Planning a grocery budget matters when studying or starting work. In this hypothetical setting, a fixed register lists all 280,000 young adults registered as living in Medellin throughout September 2026. For each person, record total September grocery spending in thousand COP. Separately, record whether the person belonged to a grocery rewards program in September: yes or no.

The initial independent random sample contains 196 people. Their mean grocery spending is 310 thousand COP, with sample standard deviation 92 thousand COP; 49 belonged to the rewards program. To compare results, draw a second independent random sample from the same register for the same September period. It also contains 196 people, with mean grocery spending 326 thousand COP, sample standard deviation 103 thousand COP and 98 rewards-program members.

At the start, only a certified report about population spread is available. It makes the population standard deviation of September grocery spending known as 98 thousand COP; it does not disclose the population mean. The two sample standard deviations remain separate descriptive quantities. The initial study estimates the initially unknown mean at 80\% and 98\% confidence.

Assume the grocery-spending distribution is sufficiently regular, without dominant outliers, for a normal approximation to sample means of this size. Each study draws a simple random sample without replacement from the full fixed register. Samples are selected independently across studies; overlap between different samples is permitted. Each study selects only 0.07\% of the population, so dependence within a draw and the finite-population correction are negligible.

Ten subsequent independent studies each use 196 people, the same known population standard deviation and both the 80\% and 98\% normal confidence procedures. These are additional studies, separate from both samples above. Each study builds both intervals around its own sample mean. The 98\% interval is wider and contains that study's 80\% interval. A later comprehensive investigation of the entire September population establishes its mean grocery spending as 315 thousand COP. Repeated samples do not establish this value. The table lists both confidence intervals for each of these ten subsequent studies, in study order. All interval endpoints are in thousand COP.

\examtable{@{}lKK@{}}{
\textbf{Study \#}&\textbf{98\% confidence interval}&\textbf{80\% confidence interval}\\\midrule
1&[283.718, 316.282]&[291.026, 308.974]\\
2&[295.718, 328.282]&[303.026, 320.974]\\
3&[303.718, 336.282]&[311.026, 328.974]\\
4&[317.718, 350.282]&[325.026, 342.974]\\
5&[285.718, 318.282]&[293.026, 310.974]\\
6&[301.718, 334.282]&[309.026, 326.974]\\
7&[291.718, 324.282]&[299.026, 316.974]\\
8&[313.718, 346.282]&[321.026, 338.974]\\
9&[281.718, 314.282]&[289.026, 306.974]\\
10&[299.718, 332.282]&[307.026, 324.974]\\
}

Complete all five \textbf{General tasks for every problem} above.
\endgroup

\columnbreak

\SubsectionBox{Young adults' phone costs \score{C1-\sml{4}\quad C2-\sml{4}}}
\QuestionLabel{q:2}

A mobile provider wants to understand phone expenses for university students and young employees. Its Cartagena branch network administers all 800,000 active personal mobile plans held by these customers in this hypothetical register on September 30, 2026. The amount charged to each plan during September 2026, in thousand COP, is the numerical measurement. A separate yes-or-no entry shows whether automatic renewal was enabled on September 30.

For the initial independent random sample, select 400 eligible plans. The mean September charge is 110 thousand COP and the sample standard deviation is 74 thousand COP; 120 plans had automatic renewal enabled. Repeat the selection independently for a second independent random sample of 400 plans from that same fixed register and period. This sample gives mean charge 105 thousand COP, sample standard deviation 83 thousand COP and 260 plans with automatic renewal enabled.

The starting information comes from a certified report about spread only. It establishes the population standard deviation of September plan charges as 80 thousand COP while withholding the population mean. This known spread differs from both sample standard deviations. The initial study estimates the initially unknown population mean at 80\% and 98\% confidence.

Assume a sufficiently regular plan-charge distribution without dominant outliers and a valid normal approximation for sample means of this size. Each study draws a simple random sample without replacement from the same full fixed register; draws across studies are independent and may overlap. Each sample contains only 0.05\% of the population, so the finite-population correction is negligible.

Two hundred fifty subsequent independent studies, separate from the two samples above, each select 400 plans and use the same known population standard deviation and both confidence procedures. Each study builds both intervals around its own sample mean. The 98\% interval is wider and contains that study's 80\% interval. A later comprehensive investigation of all eligible plans establishes mean September plan charges as 113 thousand COP; the repeated samples do not establish it. The table records mutually exclusive retrospective coverage categories for these same 250 studies.

\examtable{Kr}{
\textbf{Intervals covering the later established mean}&\textbf{Studies}\\\midrule
Both confidence levels&185\\
Only the 98\% level&55\\
Neither confidence level&10\\}

Complete all five \textbf{General tasks for every problem} above.

\columnbreak

\SubsectionBox{Create an integrated scenario from scratch\newline \score{C1-\sml{6}\quad C2-\sml{6}}}
\QuestionLabel{q:3}

Invent a new hypothetical economic or financial scenario integrating both criteria. First supply a complete, valid economic or financial setting, including the shared contextual description and all data required by the checklist. Then complete all five shared general tasks for your own scenario. One confidence level is sufficient. Write the contextual description in ordinary language: numerals, percentages, written units and plain-language table headings are allowed, but reserve symbols, formulas, inequalities and interval notation for task reference material and solutions.

\textbf{Required construction checklist}
\begin{enumerate}
\item Define a population's place, fixed period and size; define one numerical measurement with units and a separate binary condition.
\item Provide an initial independent random sample and a second independent random sample from exactly the same population and period. For each give size, mean, sample standard deviation and integer condition count between zero and that sample size. Ensure summaries are feasible for the numerical measurement.
\item State a known population standard deviation, distinguish it from both sample standard deviations, and explain that the population mean is initially unknown. Choose confidence and a supplied standard normal critical value.
\item Give suitable randomness, independence and normal-approximation assumptions. For finite sampling without replacement, keep the sampling fractions negligible. Subsequent studies must be independent draws from the same full fixed population and period.
\item Specify a number of subsequent studies, each with the initial sample size, known population standard deviation and confidence procedure. Supply all subsequent endpoints or complete coverage/miss counts. Endpoints must have the stated procedure's width; with multiple levels, evidence must respect nesting.
\item Include a later comprehensive population investigation establishing the mean of that same population for retrospective coverage. A clearly hypothetical investigation is acceptable; do not claim repeated samples establish the population mean.
\end{enumerate}

\par\smallskip
Complete all five \textbf{General tasks for every problem} above.




\end{multicols}
\end{document}
